% !TEX root = ../main.tex
%----------------------------------------------------------
% 04-pignistic-bridge.tex — Section 3: Pignistic Bridge
% MIGRATED from preprint A3-information-theory.tex
% Removed: "Product 2"/"Product 3" framing, Clyfar-specific references
% Ozone categories → SPC categories
%----------------------------------------------------------

\section{Pignistic Bridge}\label{sec:pignistic_bridge}

The term \textit{pignistic} derives from the Latin \textit{pignus} (a
wager): it denotes the probabilities one would assign if forced to bet.
Standard verification scores---the logarithmic score, Brier score, and
information gain---require probability inputs that sum to unity.  A
possibility distribution does not satisfy this constraint: its values
encode plausibility rather than relative frequency, and a subnormal
distribution does not even reach a maximum of one.  The
\textit{possibility-to-probability bridge} developed below converts a
subnormal possibility distribution into an $(n{+}1)$-category probability
vector, preserving the subnormality signal as an explicit ignorance
outcome rather than erasing it through na\"ive normalisation.

%----------------------------------------------------------
\subsection{Possibility-to-Probability Bridge}\label{sec:poss_prob_bridge}

Applying probabilistic verification to possibilistic forecasts requires
converting possibility distributions to probabilities.  The
\textit{possibility-to-probability bridge} preserves the subnormality
signal through three steps:

\begin{enumerate}[leftmargin=2em]
\item \textbf{Reserve ignorance mass:}
    $p_{\mathrm{ign}} = \ign = 1 - \max(\pi)$.
\item \textbf{Distribute the remaining $(1 - \ign)$ proportionally
    across outcomes:}
    \[
    p_i = \frac{\pi_i \cdot (1 - \ign)}{\sum_j \pi_j}
    \]
\item \textbf{Append $p_{\mathrm{ign}}$ as an explicit ``ignorance
    outcome,''} yielding a probability distribution over $n+1$
    categories that sums to unity.
\end{enumerate}

Simple normalisation ($p_i = \pi_i / \sum_j \pi_j$) erases
subnormality---probabilities always sum to one regardless of model
confidence.  This is equivalent to the pignistic transformation of
\citet{Smets1990-MISSING}, which was designed for normal belief
functions.  The bridge preserves the ignorance signal,
ensuring that verification scores correctly penalise uncertain forecasts
by routing probability mass into a non-observed outcome.  For normal
distributions ($\max \pi = 1$, $\ign = 0$), both methods coincide.

% [PLACEHOLDER: 1–2 sentences on the connection to transferable belief
%  model (Smets & Kennes 1994) and why our extension to subnormal
%  distributions is novel.]

\textbf{Worked example.}
Suppose a possibilistic convective outlook produces
$\pi = (0.2,\; 0.4,\; 0.6,\; 0.1,\; 0.0)$.
Then:
\begin{itemize}[leftmargin=2em]
\item $\ign = 1 - 0.6 = 0.4$ (system is 40\% unsure).
\item $\sum_j \pi_j = 1.3$; remaining $0.6$ distributed
    proportionally:
\item $\mathbf{p} = (0.092,\; 0.185,\; 0.277,\; 0.046,\; 0.0,\;
    \underbrace{0.400}_{\text{ign}})$.
\item If \spc{ENH} is observed: self-entropy
    $= -\log_2(0.277) = 1.85$~bits.
\item Climatological baseline (e.g.,
    $p_{\mathrm{clim}}(\spc{ENH}) = 0.05$):
    $-\log_2(0.05) = 4.32$~bits.
\item $\mathrm{IG} = 4.32 - 1.85 = 2.47$ bits gained over
    climatology.
\end{itemize}

High \ign{} ($0.4$) correctly inflated the ignorance outcome, reducing
$p(\spc{ENH})$ from what simple normalisation would give
($0.6/1.3 = 0.462$) to $0.277$.  The verification score ($1.85$~bits)
honestly reflects the system's partial uncertainty.
Figure~\ref{fig:pignistic_bridge} visualises this transformation.

\begin{figure*}[t]
\centering
\includegraphics[width=0.85\textwidth]{fig4_pignistic_bridge.png}
\caption{Possibility-to-probability bridge applied to a subnormal possibility
    distribution.  \textbf{Left:} raw possibility values $\pi(\omega)$
    for five SPC categories.  \textbf{Right:} resulting $(n{+}1)$-category
    probability vector, with the grey bar representing the explicit
    ignorance outcome $p_{\mathrm{ign}} = \ign$.  Arrows show
    proportional redistribution of the remaining $(1 - \ign)$ mass.  The
    ignorance bar absorbs probability mass that simple normalisation would
    have spread across all categories, preserving the subnormality signal.}
\label{fig:pignistic_bridge}
\end{figure*}

%----------------------------------------------------------
\subsection{Information-Gain Decomposition}\label{sec:ig_decomposition}

% [PLACEHOLDER: Opening paragraph — briefly introduce the information-gain
%  framework and its relationship to the logarithmic scoring rule.]

After the bridge converts possibilities to probabilities,
standard information-theoretic verification applies.  Information gain
measures the reduction in surprise from replacing one forecast with
another:
\begin{equation}
\mathrm{IG} = D_{\mathrm{KL}}(o \,\|\, f_1) - D_{\mathrm{KL}}(o \,\|\, f_2)
\label{eq:ig}
\end{equation}
where $o$ is the observation, $f_1$ is the baseline (e.g., climatology),
and $f_2$ is the possibilistic forecast (after bridge conversion).
Positive IG indicates that the forecast reduced surprise relative to the
baseline.
% [PLACEHOLDER: Add Roulston & Smith (2002) and Lawson (2024) citations.]

Information gain decomposes diagnostically:
\begin{equation}
\mathrm{IG}_{\mathrm{decomp}} = \mathrm{UNC} - \mathrm{DSC} + \mathrm{REL}
\quad \text{(all in bits)}
\label{eq:ig_decomp}
\end{equation}
where:
\begin{itemize}[leftmargin=2em]
\item \textbf{UNC} is irreducible uncertainty (entropy of the base rate).
\item \textbf{REL} is calibration error (lower is better): are stated
    probabilities honest?
\item \textbf{DSC} is discrimination (higher is better): can the system
    tell events apart?
\end{itemize}

This decomposition diagnoses whether a possibilistic forecast's value
arises from calibration or from discriminating between severe-weather
regimes.  Figure~\ref{fig:ig_decomp} illustrates the decomposition for
several forecast scenarios with net IG annotated.

% [PLACEHOLDER: Brief paragraph connecting IG decomposition to familiar
%  meteorological analogues — Brier decomposition, reliability diagrams.]

\begin{figure*}[t]
\centering
\includegraphics[width=0.75\textwidth]{fig7_ig_decomposition.png}
\caption{Information-gain decomposition for five forecast scenarios.
    Stacked bars show the three components: UNC (light grey;
    irreducible uncertainty), DSC (purple; discrimination---higher is
    better), and REL (green; calibration error---lower is better).  Net
    IG is annotated above each bar.  The sharp-correct scenario achieves
    the highest IG; the sharp-wrong scenario shows negative IG (worse
    than climatology).}
\label{fig:ig_decomp}
\end{figure*}

To avoid naming collision with possibilistic ignorance $\ign$, this
manuscript uses ``information gain'' (IG) rather than the abbreviation
IGN used in some verification literature.

%----------------------------------------------------------
\subsection{Why Na\"ive Normalisation Erases Information}%
\label{sec:naive_normalisation}

Simple normalisation ($p_i = \pi_i / \sum_j \pi_j$) is the standard
route from possibilities to probabilities.  It forces the result to sum
to one regardless of model confidence: a system that is $90\%$ unsure
($\ign = 0.9$) produces the same probability vector as one that is
$10\%$ unsure ($\ign = 0.1$), provided their raw shapes are
proportional.  The ignorance signal is erased.  Verification scores
applied to the normalised vector cannot distinguish a sharp, confident
forecast from a diffuse, hedged one---the answer to ``how would you
bet?'' is given without acknowledging ``I wouldn't bet at all.''

The possibility-to-probability bridge avoids this by deferring
probabilistic commitment for the ignorance portion.  Mass equal to
$\ign$ is routed into an $(n{+}1)$-th outcome that can never be the
observed category.  Uncertain forecasts are penalised because this
stranded mass earns no verification credit: it inflates the self-entropy
of the converted forecast in proportion to $\ign$.  The more ignorant
the system, the more mass is stranded, and the larger the penalty.

\textbf{Worked contrast.}
Using $\pi = (0.2,\, 0.4,\, 0.6,\, 0.1,\, 0.0)$ from the worked
example in Section~\ref{sec:poss_prob_bridge}:
\begin{itemize}[leftmargin=2em]
\item Simple normalisation: $p(\spc{ENH}) = 0.6 / 1.3 = 0.462$.
\item Bridge: $p(\spc{ENH}) = 0.277$, $p_{\mathrm{ign}} = 0.400$.
\item If \spc{ENH} is observed: self-entropy $= -\log_2(0.462) = 1.11$~bits
    (simple) vs.\ $-\log_2(0.277) = 1.85$~bits (bridge).
\item The $0.74$-bit gap is the verification cost of ignoring the
    system's admitted uncertainty.  Simple normalisation rewards the
    system for confidence it never claimed.
\end{itemize}

\textbf{Complementarity of \ign{} and IG.}
Possibilistic ignorance $\ign$ and information gain IG capture
complementary dimensions of forecast quality:
\begin{itemize}[leftmargin=2em]
\item \textbf{$\ign$ (possibilistic ignorance):} Measures the
    forecasting system's support strength---did the system know what was
    going on?  Derived from the subnormality of the raw output
    (Eq.~\ref{eq:ignorance}).
\item \textbf{IG (information gain):} Measures forecast value---did the
    forecast reduce surprise about what happened?  Derived from the
    probability-converted forecast evaluated against observations.
\end{itemize}

Both are needed because:
\begin{itemize}[leftmargin=2em]
\item A \textbf{sharp wrong} forecast has low \ign{} but negative IG.
\item A \textbf{hedged correct} forecast has high \ign{} but positive IG.
\item A \textbf{sharp correct} forecast has low \ign{} and high IG---the
    ideal.
\item A \textbf{hedged wrong} forecast has high \ign{} and negative
    IG---the system knew it didn't know, and it still got it wrong.
\end{itemize}
